\documentclass[12pt,spanish,letterpaper,oneside]{article}
\def\documenttitle{Examen Parcial 2:\\[6pt]Evaluaci\'on y fundamentaci\'on dogm\'atica}
\def\documentsubtitle{Fundamentaci\'on dogm\'atica y orden jur\'idico}
\def\documenttitlehf{Examen Parcial 2: Derechos Humanos}
\def\hfwidthwrap{false}
\def\hfwidthtitle{0.50}
\def\hfwidthcourse{0.45}
\def\documentsubject{Licenciatura en Derecho}
\def\documentauthor{Mart\'in Jonathan de la Cruz Mu\~noz}
\def\documentauthorcanonico{Martín Jonathan de la Cruz Muñoz}
\author{Mart\'in Jonathan de la Cruz Mu\~noz}
\title{Examen Parcial 2: Evaluaci\'on y fundamentaci\'on dogm\'atica}
\def\coursename{Antecedentes de los derechos humanos}
\def\coursecode{005}
\def\universityname{Universidad Abierta y a Distancia de M\'exico}
\def\universityfaculty{Licenciatura en Derecho}
\def\universitydepartment{Antecedentes de los derechos humanos}
\def\universitydepartmentimage{departamentos/UnADM}
\def\universitydepartmentimagecfg{height=1.57cm}
\def\universitylocation{Ciudad de M\'exico}
\def\authortable{\begin{tabular}{ll}
\textbf{Alumno:} & Mart\'in Jonathan de la Cruz Mu\~noz \\
\textbf{Matr\'icula:} & ES2611202040 \\
\textbf{Figura docente:} & Daniel Alexis Lozano Ortega \\
\textbf{Grupo:} & 005 \\
\textbf{Semestre / Bloque:} & 2 / 2 \\
\textbf{Semana lectiva:} & Semana 7 \\
\textbf{Fecha de realizaci\'on:} & Noviembre de 2026 \\
\multicolumn{2}{l}{\universitylocation}
\end{tabular}}

\usepackage{iftex}
\ifLuaTeX
	\let\pdfpagewidth\pagewidth
	\let\pdfpageheight\pageheight
\fi

% Wrapper local para apuntar a la plantilla maestra de AulaTeX
\input{C:/Users/Sysx/Documents/AulaTeX-Academico/aulaTEX/base/Plantilla-Informe/template.tex}


% Refuerzo post-template
\def\documentsubtitle{Fundamentaci\'on dogm\'atica y orden jur\'idico}
\def\hfwidthwrap{false}
\def\hfwidthtitle{0.50}
\def\hfwidthcourse{0.45}

\usepackage{booktabs}
\usepackage{tabularx}
\usepackage{array}
\usepackage{colortbl}
\usepackage{xcolor}

\definecolor{azulOscuro}{RGB}{20, 45, 85}
\definecolor{azulMedio}{RGB}{40, 90, 150}
\definecolor{grisFondo}{RGB}{245, 247, 250}
\definecolor{tableHeader}{RGB}{232, 238, 246}
\definecolor{bordeGris}{RGB}{180, 190, 205}

\begin{document}
\templatePortrait
\templatePagecfg
\templateFinalcfg

% =============================================================================
% RESUMEN (ABSTRACT)
% =============================================================================
\begin{abstractd}
El presente reporte acad\'emico sistematiza y fundamenta jur\'idicamente la resoluci\'on del Examen Parcial 2 correspondiente a la asignatura Antecedentes de los derechos humanos en la Universidad Abierta y a Distancia de M\'exico (UnADM). La evaluaci\'on integra los contenidos estructurales abordados en las semanas 5, 6 y 7 del plan formativo, examinando el orden jur\'idico nacional, los mecanismos de garant\'ia, la evoluci\'on hist\'orica del reconocimiento de derechos y los par\'ametros de protecci\'on internacional e interamericana. A trav\'es de un an\'alisis dogm\'atico estructurado en tres tiempos anal\'iticos, se desglosan los cinco reactivos que integran la evaluaci\'on oficial, justificando cada respuesta correcta frente al marco constitucional mexicano derivado de la reforma de junio de 2011, la dignidad de la persona como norma jur\'idica vinculante y el principio de complementariedad convencional. El documento concluye con una matriz integral de resultados y una ponderaci\'on cr\'itica rigurosa sobre la funci\'on del garantismo jurisdiccional y el principio de \emph{ultima ratio} como baluartes indispensables frente a las tentaciones punitivas del poder estatal.
\end{abstractd}

% =============================================================================
% ÍNDICE DE CONTENIDOS
% =============================================================================
\templateIndex

% =============================================================================
% TIEMPO 1: INTRODUCCIÓN Y FUNDAMENTACIÓN DOGMÁTICA
% =============================================================================
\clearpage
\section[Introducci\'on y fundamentaci\'on]{Introducci\'on y fundamentaci\'on dogm\'atica de la evaluaci\'on formativa}

La evaluaci\'on en la ense\~nanza universitaria del Derecho contempor\'aneo no puede concebirse como un mero tr\'amite nemot\'ecnico o una recopilaci\'on mec\'anica de respuestas automatizadas. En el modelo formativo de la Universidad Abierta y a Distancia de M\'exico, el examen parcial representa una instancia procesal de autoverificaci\'on cr\'itica que permite constatar la asimilaci\'on epistemol\'ogica y hermen\'eutica de los contenidos desarrollados a lo largo de las semanas 5, 6 y 7 del curso. Sin embargo, dicha evaluaci\'on carecer\'ia de sentido sin un an\'alisis anal\'itico profundo; por ello, este reporte articula la transici\'on fundamental entre el sustento filos\'ofico-hist\'orico de las libertades p\'ublicas y su concreci\'on procesal en el orden constitucional mexicano y los tratados internacionales de los que el Estado es parte \citep{alvarezIcaza2009ADH}. Dado que el programa de estudios exige vincular teor\'ia y pr\'actica jur\'idica, la ponderaci\'on dogm\'atica de cada reactivo evidencia la solidez del razonamiento empleado.

El Examen Parcial 2 concentra su radio anal\'itico en tres n\'ucleos tem\'aticos de inexcusable dominio para el jurista: primero, la g\'enesis y universalizaci\'on de los derechos a trav\'es de los grandes hitos emancipadores y documentos fundacionales internacionales; segundo, la naturaleza ontol\'ogica de los derechos como prerrogativas inherentes a la dignidad humana, innegociables e irreductibles a consideraciones de car\'acter patrimonial o contractual; y tercero, la conformaci\'on del bloque de regularidad constitucional en M\'exico a partir de la trascendental redacci\'on del art\'iculo 1.º de la Constituci\'on Pol\'itica de los Estados Unidos Mexicanos (CPEUM), articulado arm\'onicamente con los \'organos rectores del Sistema Interamericano de Derechos Humanos (SIDH) \citep{carbonell2019, ferrajoli2001}.

Las respuestas analizadas en este documento fueron sustentadas y registradas formalmente en el entorno virtual de aprendizaje de la UnADM dentro de los plazos reglamentarios. No obstante las restricciones de tiempo inherentes a las plataformas digitales, el objetivo de este informe reside en superar la mera selecci\'on de opciones m\'ultiples mediante la construcci\'on de una s\'olida fundamentaci\'on dogm\'atica y jurisprudencial para cada reactivo. De esta forma, se evidencia que la selecci\'on de cada respuesta deriva de una interpretaci\'on sistem\'atica de las fuentes doctrinales y normativas vigentes, vinculando la evaluaci\'on formativa con la praxis forense y la estricta vigencia del principio de legalidad \citep{cpeumADH2026}.

% =============================================================================
% TIEMPO 2: DESARROLLO ORGÁNICO Y RESOLUCIÓN DOGMÁTICA DE LOS REACTIVOS
% =============================================================================
\section[Desarrollo org\'anico]{Desarrollo org\'anico y resoluci\'on dogm\'atica de los reactivos del Examen Parcial 2}

El examen oficial const\'o de cinco reactivos (cuatro de opci\'on m\'ultiple y uno de respuesta escrita breve) dise\~nados para contrastar categor\'ias normativas precisas. A continuaci\'on, se presenta el an\'alisis individualizado de cada reactivo, precisando la formulaci\'on institucional, las opciones disponibles, la determinaci\'on jur\'idica de la respuesta correcta, el nivel de confianza y la fundamentaci\'on te\'orica exhaustiva.

\subsection{Reactivo 1: Documentos fundamentales y g\'enesis moderna de los derechos}
\textbf{Pregunta institucional:} ¿Cu\'al de los siguientes documentos es considerado un antecedente fundamental para el desarrollo moderno de los derechos humanos?
\begin{enumerate}
\item[a)] Ley Federal del Trabajo.
\item[b)] C\'odigo Civil Federal.
\item[c)] C\'odigo Penal Federal.
\item[d)] Declaraci\'on Universal de los Derechos Humanos.
\end{enumerate}

\textbf{Respuesta correcta:} \textbf{d) Declaraci\'on Universal de los Derechos Humanos.} \\
\textbf{Nivel de confianza:} Alta.

\textbf{Fundamentaci\'on dogm\'atica:} Proclamada el 10 de diciembre de 1948 por la Asamblea General de las Naciones Unidas mediante la Resoluci\'on 217 A (III), la Declaraci\'on Universal de los Derechos Humanos (DUDH) constituye el documento axial e insustituible sobre el cual se edific\'o el derecho internacional contempor\'aneo tras los horrores de la Segunda Guerra Mundial \citep{onuDeclaracion1948ADH}. Las opciones restantes representan ordenamientos adjetivos y sustantivos ordinarios de derecho interno en M\'exico (laboral, civil y punitivo), cuyas normas regulan esferas sectoriales pero carecen del car\'acter de matriz fundacional universal de los derechos modernos. En consecuencia, la DUDH transform\'o el estatus jur\'idico del individuo frente a la soberan\'ia estatal al consagrar un cat\'alogo universal indivisible e interdependiente.

\subsection{Reactivo 2: La inherencia ontol\'ogica de los derechos y la dignidad humana}
\textbf{Pregunta institucional:} ¿Cu\'al de las siguientes caracter\'isticas corresponde a los derechos humanos?
\begin{enumerate}
\item[a)] Dependen de la condici\'on econ\'omica de cada persona.
\item[b)] Son inherentes a todas las personas.
\item[c)] Solo pueden reconocerse mediante contratos.
\item[d)] Son exclusivos de determinadas personas.
\end{enumerate}

\textbf{Respuesta correcta:} \textbf{b) Son inherentes a todas las personas.} \\
\textbf{Nivel de confianza:} Alta.

\textbf{Fundamentaci\'on dogm\'atica:} La inherencia es la propiedad ontol\'ogica consustancial de los derechos humanos. Como ha sostenido Luigi Ferrajoli en su garantismo jur\'idico, los derechos fundamentales son aquellos cuya titularidad corresponde universalmente a todos los seres humanos en cuanto personas dotadas del estatus de personas con capacidad de obrar \citep{ferrajoli2001}. Es decir, la Primera Sala de la Suprema Corte de Justicia de la Naci\'on consolid\'o este criterio en la jurisprudencia 1a./J. 37/2016 (10a.), declarando que la dignidad humana no es una simple declaraci\'on \'etica o axiol\'ogica, sino una norma jur\'idica vinculante que consagra un derecho fundamental del cual emanan todas las dem\'as prerrogativas b\'asicas \citep{scjnDignidad2016ADH}. Se descarta radicalmente que dependan del patrimonio (opci\'on a), de pactos sinalagm\'aticos bilaterales (opci\'on c) o de privilegios estamentales excluyentes (opci\'on d).

\subsection{Reactivo 3: El art\'iculo 1.º constitucional como par\'ametro de regularidad}
\textbf{Pregunta institucional:} El reconocimiento de los derechos humanos y sus principios en la Constituci\'on Pol\'itica de los Estados Unidos Mexicanos se encuentra establecido en el art\'iculo \underline{\hspace{1.5cm}}º. (s\'olo debe escribir el n\'umero ar\'abigo, no en letra)
\begin{enumerate}
\item[4.] 1
\item[5.] 2
\item[6.] 3
\item[7.] 4
\end{enumerate}

\textbf{Respuesta correcta:} \textbf{1} (escrita en d\'igito ar\'abigo simple: \textbf{1}). \\
\textbf{Nivel de confianza:} Alta.

\textbf{Fundamentaci\'on dogm\'atica:} El art\'iculo 1.º de la CPEUM, profundamente reformado el 10 de junio de 2011, representa el epicentro dogm\'atico de la justicia constitucional mexicana. Dicho precepto estatuye que en los Estados Unidos Mexicanos todas las personas gozar\'an de los derechos reconocidos en la propia Constituci\'on y en los tratados internacionales de los que el Estado sea parte, as\'i como de las garant\'ias para su protecci\'on. Asimismo, dicho art\'iculo positiviza el mandato hermen\'eutico de interpretaci\'on conforme, el principio \emph{pro persona} y las obligaciones expresas para todas las autoridades de promover, respetar, proteger y garantizar los derechos bajo los principios de universalidad, interdependencia, indivisibilidad y progresividad \citep{cpeumADH2026, scjnParametro2014ADH}. La indicaci\'on requiri\'o expresamente consignar el d\'igito num\'erico ar\'abigo ``1'', respetando el formato autom\'atico de calificaci\'on del entorno Moodle.

\subsection{Reactivo 4: La Revoluci\'on Francesa y las libertades del constitucionalismo moderno}
\textbf{Pregunta institucional:} ¿Cu\'al de los siguientes acontecimientos hist\'oricos contribuy\'o al desarrollo de las ideas modernas sobre los derechos y libertades de las personas?
\begin{enumerate}
\item[a)] Creaci\'on de la Organizaci\'on Mundial de la Salud.
\item[b)] Revoluci\'on Francesa.
\item[c)] Creaci\'on de las redes sociales.
\item[d)] Invenci\'on de internet.
\end{enumerate}

\textbf{Respuesta correcta:} \textbf{b) Revoluci\'on Francesa.} \\
\textbf{Nivel de confianza:} Alta.

\textbf{Fundamentaci\'on dogm\'atica:} La Revoluci\'on Francesa de 1789 signific\'o la ruptura definitiva con el absolutismo mon\'arquico y el feudalismo corporativo en Europa continental. La adopci\'on de la Declaraci\'on de los Derechos del Hombre y del Ciudadano del 26 de agosto de 1789 proclam\'o la igualdad ante la ley, la libertad individual, la presunci\'on de inocencia y la soberan\'ia nacional \citep{declaracionFrancia1789ADH}. Su c\'elebre art\'iculo 16 condens\'o la f\'ormula indeleble del constitucionalismo occidental: toda sociedad donde no est\'e garantizada la separaci\'on de poderes ni asegurada la protecci\'on de los derechos carece en rigor de Constituci\'on. Las dem\'as opciones constituyen desarrollos tecnol\'ogicos de finales del siglo XX o agencias multilaterales especializadas posteriores a la Segunda Guerra Mundial (OMS en 1948), que no representan el acontecimiento pol\'itico-filos\'ofico fundacional del constitucionalismo moderno de derechos.

\subsection{Reactivo 5: \'Organos del Sistema Interamericano de Derechos Humanos}
\textbf{Pregunta institucional:} ¿Cu\'al de los siguientes organismos pertenece al Sistema Interamericano de Derechos Humanos?
\begin{enumerate}
\item[a)] Organizaci\'on Internacional del Trabajo.
\item[b)] Organizaci\'on Mundial del Comercio.
\item[c)] Comisi\'on Interamericana de Derechos Humanos.
\item[d)] Banco Mundial.
\end{enumerate}

\textbf{Respuesta correcta:} \textbf{c) Comisi\'on Interamericana de Derechos Humanos.} \\
\textbf{Nivel de confianza:} Alta.

\textbf{Fundamentaci\'on dogm\'atica:} El Sistema Interamericano de Derechos Humanos (SIDH), erigido en el seno de la Organizaci\'on de los Estados Americanos (OEA) a partir de la Carta de la OEA y la Convenci\'on Americana sobre Derechos Humanos (Pacto de San Jos\'e de 1969), cuenta con dos \'organos cardinales con competencias especializadas: la Comisi\'on Interamericana de Derechos Humanos (CIDH, con sede en Washington, D.C.) y la Corte Interamericana de Derechos Humanos (Corte IDH, con sede en San Jos\'e, Costa Rica) \citep{oeaPacto1969ADH}. La CIDH tramita peticiones individuales, emite medidas cautelares y redacta informes sobre la situaci\'on de derechos en los Estados miembros. En contraste, la OIT pertenece al sistema especializado de la ONU, mientras que la OMC y el Banco Mundial son organismos financieros y comerciales internacionales ajenos al esquema convencional interamericano.

% -----------------------------------------------------------------------------
% PRODUCTO VISIBLE INTEGRADO: MATRIZ DE RESOLUCIÓN Y SÍNTESIS
% -----------------------------------------------------------------------------
\vspace{8pt}
\begin{table}[htbp]
\centering
\caption{Matriz sint\'etica de resoluci\'on y fundamentaci\'on del Examen Parcial 2}\label{tab:matrizExamen}
\vspace{3pt}
\begingroup
\scriptsize
\setlength{\tabcolsep}{3.5pt}
\renewcommand{\arraystretch}{1.20}
\begin{tabularx}{\linewidth}{>{\bfseries\centering\arraybackslash}p{0.7cm} >{\raggedright\arraybackslash}p{2.3cm} >{\raggedright\arraybackslash}p{2.9cm} >{\centering\arraybackslash}p{1.3cm} >{\raggedright\arraybackslash}X}
\toprule
\rowcolor{tableHeader}
No. & Dimensi\'on tem\'atica & Opci\'on / Concepto clave & Confianza & Base dogm\'atica y fuente aplicable \\
\midrule
1 & Antecedentes modernos universales & Declaraci\'on Universal de los Derechos Humanos (1948) & Alta & Resoluci\'on 217 A (III) de la Asamblea General de la ONU; g\'enesis del derecho internacional humanitario y convencional. \\
\addlinespace
2 & Naturaleza ontol\'ogica de los derechos & Inherentes a todas las personas & Alta & Dignidad humana como norma jur\'idica vinculante; tesis 1a./J. 37/2016 y teor\'ia garantista de Ferrajoli. \\
\addlinespace
3 & Par\'ametro de regularidad constitucional & Art\'iculo 1.º de la CPEUM (D\'igito ar\'abigo: 1) & Alta & Bloque constitucional mexicano reformado en 2011; principio \emph{pro persona} y jurisprudencia P./J. 20/2014. \\
\addlinespace
4 & Constitucionalismo y revoluciones & Revoluci\'on Francesa (1789) & Alta & Declaraci\'on de los Derechos del Hombre y del Ciudadano; principio de legalidad, igualdad y separaci\'on de poderes. \\
\addlinespace
5 & Sistema Interamericano de Derechos & Comisi\'on Interamericana de Derechos Humanos (CIDH) & Alta & Art\'iculos 33 a 51 de la Convenci\'on Americana sobre Derechos Humanos (Pacto de San Jos\'e de 1969). \\
\bottomrule
\end{tabularx}
\endgroup
\par\vspace{3pt}
{\footnotesize\raggedright\textbf{Fuente:} Elaboraci\'on propia a partir de las consignas oficiales del Examen Parcial 2 de UnADM y el contraste normativo correspondiente.\par}
\end{table}

% =============================================================================
% TIEMPO 3: CONCLUSIONES Y PONDERACIÓN CRÍTICA
% =============================================================================
\clearpage
\section[Conclusiones]{Conclusiones y ponderaci\'on cr\'itica sobre la evaluaci\'on formativa y el garantismo}

La resoluci\'on del Examen Parcial 2 confirma que el aprendizaje del Derecho de los derechos humanos no consiste en asimilar dogmas est\'aticos ni en memorizar mec\'anicamente cat\'alogos normativos desarticulados de su realidad social. En consecuencia, sostengo con firmeza que en el ejercicio profesional y en el foro judicial la comprensi\'on genuina de estos principios exige articular la evoluci\'on hist\'orica de las luchas emancipadoras con las herramientas procesales efectivas de protecci\'on constitucional. Una evaluaci\'on formativa adquiere verdadero sentido cuando trasciende la calificaci\'on cuantitativa de plataforma y estimula un juicio cr\'itico fundamentado sobre los l\'imites que el ordenamiento jur\'idico impone al poder coercitivo del Estado.

Asimismo, la vinculaci\'on intr\'inseca entre los \'organos del Sistema Interamericano y el art\'iculo 1.º constitucional demuestra que M\'exico no puede operar como un enclave herm\'etico ajeno a los est\'andares internacionales de justicia. Considero insostenible cualquier intento de relativizar la dignidad de la persona so pretexto de doctrinas de seguridad nacional o emergencias pol\'iticas coyunturales. La dignidad humana es una norma jur\'idica imperativa y operativa que proscribe tratar al individuo como un instrumento para fines colectivos; su efectividad demanda una judicatura valiente, capaz de ejercer un control difuso de convencionalidad riguroso y coherente con las l\'ineas jurisprudenciales interamericanas.

Por tanto, rechazo frontalmente la instrumentalizaci\'on del discurso penal como supuesta respuesta autom\'atica a los reclamos sociales, as\'i como cualquier pr\'actica que relaje las garant\'ias judiciales bajo la falacia de una mayor eficacia punitiva. La estricta legalidad y el principio de \emph{ultima ratio} constituyen las mayores conquistas del pensamiento civilizatorio frente a la arbitrariedad; permitir la prisi\'on preventiva oficiosa o debilitar la presunci\'on de inocencia degrada el pacto constitucional y lesiona de ra\'iz la inherencia ontol\'ogica que juramos defender.\footnote{Declaraci\'on sobre el uso de herramientas de inteligencia artificial: En cumplimiento estricto con los Lineamientos para el uso \'etico y responsable de la inteligencia artificial de la Universidad Abierta y a Distancia de M\'exico (UnADM), se hace constar que se utiliz\'o un modelo de lenguaje artificial exclusivamente como apoyo instrumental en la correcci\'on de estilo y composici\'on editorial. La autor\'ia conceptual, el an\'alisis jur\'idico-dogm\'atico, la postura cr\'itica y las conclusiones son de exclusiva responsabilidad intelectual del estudiante.}

% =============================================================================
% REFERENCIAS BIBLIOGRÁFICAS
% =============================================================================
\clearpage
\begin{thebibliography}{99}

\bibitem[Alvarez Icaza Longoria(2009)]{alvarezIcaza2009ADH}
Alvarez Icaza Longoria, E. (2009). \emph{Los derechos humanos en M\'exico}. M\'exico: Nostra Ediciones.

\bibitem[Asamblea Nacional Constituyente de Francia(1789)]{declaracionFrancia1789ADH}
Asamblea Nacional Constituyente de Francia. (1789). \emph{Declaraci\'on de los Derechos del Hombre y del Ciudadano de 26 de agosto de 1789}. Par\'is: Conseil constitutionnel.

\bibitem[C\'amara de Diputados(2026)]{cpeumADH2026}
C\'amara de Diputados del H. Congreso de la Uni\'on. (2026). \emph{Constituci\'on Pol\'itica de los Estados Unidos Mexicanos} (Texto vigente con reformas en el DOF hasta 2026). M\'exico: Secretar\'ia de Servicios Parlamentarios.

\bibitem[Carbonell(2019)]{carbonell2019}
Carbonell, M. (2019). \emph{Teor\'ia de los derechos fundamentales y justicia constitucional}. Madrid: Editorial Trotta.

\bibitem[Ferrajoli(2001)]{ferrajoli2001}
Ferrajoli, L. (2001). \emph{Los fundamentos de los derechos fundamentales}. Madrid: Editorial Trotta.

\bibitem[Naciones Unidas(1948)]{onuDeclaracion1948ADH}
Naciones Unidas. (1948). \emph{Declaraci\'on Universal de Derechos Humanos}. Resoluci\'on 217 A (III) de la Asamblea General. Par\'is: ONU.

\bibitem[Organizaci\'on de los Estados Americanos(1969)]{oeaPacto1969ADH}
Organizaci\'on de los Estados Americanos. (1969). \emph{Convenci\'on Americana sobre Derechos Humanos (Pacto de San Jos\'e de Costa Rica)}. San Jos\'e: OEA.

\bibitem[Suprema Corte de Justicia de la Naci\'on(2014)]{scjnParametro2014ADH}
Suprema Corte de Justicia de la Naci\'on. (2014). \emph{Jurisprudencia P./J. 20/2014 (10a.): Derechos humanos contenidos en la Constituci\'on y en los tratados internacionales. Constituyen el par\'ametro de control de regularidad constitucional, pero cuando en la Constituci\'on haya una restricci\'on expresa, se debe estar a ella}. Gaceta del Semanario Judicial de la Federaci\'on, D\'ecima \'Epoca, Libro 5, Tomo I, p. 202.

\bibitem[Suprema Corte de Justicia de la Naci\'on(2016)]{scjnDignidad2016ADH}
Suprema Corte de Justicia de la Naci\'on. (2016). \emph{Jurisprudencia 1a./J. 37/2016 (10a.): Dignidad humana. Constituye una norma jur\'idica que consagra un derecho fundamental a favor de las personas y no una simple declaraci\'on \'etica}. Gaceta del Semanario Judicial de la Federaci\'on, D\'ecima \'Epoca, Libro 33, Tomo II, p. 633.

\end{thebibliography}

\end{document}
